% /root/monodromy, 19 November 2026; independently audited by /root and /root/reduction_descent.
\subsection{Uniform structure on the fixed ordinary product surface}
\label{subsec:product-uniform-structure}
\status{Proved; parent proof, primary-source and code audits passed.
The uniform jet, incidence and union argument was also independently audited by the reduction/descent agent.}

We first record the general finite-group form of the prime-order
test used above.  We then use one further computed base degree to
obtain a statement in every degree; the latter statement does not
identify the entire group in every degree.

\begin{lemma}\label{lem:primitive-prime-transvection-test}
Let $N\ge4$ and $H\triangleleft G\subseteq\GL_N(\mathbf F_2)$.
Suppose that $H$ is nontrivial and generated by transvections.
Assume $G$ contains an element $g$ whose characteristic polynomial
is irreducible and not reciprocal, and whose order is divisible
by a prime $\ell$ with $\operatorname{ord}_{\ell}(2)=N$.
Then $H=G=\GL_N(\mathbf F_2)$.
\end{lemma}

\begin{proof}
Use the socle and multiplicity argument in
Lemma~\ref{lem:product-frobenius-group-test}.
If $H$ were reducible, $G$ would be contained in
$\GL_c(\mathbf F_2)\wr S_{N/c}$ for a proper divisor $c$ of $N$.
The order condition gives $\ell\equiv1\pmod N$, hence $\ell>N$;
and $\ell$ divides no $2^i-1$ with $i<N$.
It therefore divides neither factor of this wreath-product order.
This contradicts the order of $g$.
Thus $H$ is irreducible.  The same opened
\cite[Theorem~3.6, p.~22]{KantorSeress2001}, uniqueness of the
invariant alternating form in its non-linear alternatives, and
nonreciprocity exclude those alternatives exactly as before.
\end{proof}

\begin{corollary}\label{cor:product-degree-three-full-group}
For the fixed surface and embedding in
Theorem~\ref{thm:product-char2-two-full-groups}, the geometric and
arithmetic groups in degree three are both
$\GL_{80}(\mathbf F_2)$, over every perfect extension of
$\mathbf F_2$.
\end{corollary}

\begin{proof}
Here $W$ has dimension $81$.  The one-factor ordered basis is
\[
                 (1,x,x^2,x^3,x^4,y,xy,x^2y,x^3y).
\]
Use the same tensor and binary coefficient conventions as before.
The ordinary section encoded by $1284714998244677440413117$
has characteristic polynomial $(T+1)f_{80}(T)$, where the binary
coefficient encoding of $f_{80}$ is
\[
                         1592084315789834604401113.
\]
The retained program verifies that it is irreducible and not
reciprocal.  Its root order is divisible by
$\ell=4278255361$, and
$\operatorname{ord}_{\ell}(2)=80$.
The independent verifier checks primality by trial division
through $\lfloor\sqrt\ell\rfloor$, checks the order by the
exponents $80,40,16$, and checks the root-order assertion by
modular polynomial exponentiation.  The ordinary section
$810021528320842360421666$ has a different relative characteristic
polynomial, so the geometric image is nontrivial.

For the line with
\[
 a=998130104707720808648124,
 \qquad b_0=2032072755533295013746953,
\]
the determinant is monic of degree $81$.  Its six irreducible
factors have degrees $5,10,10,11,11,34$, and at a root of every
factor both $M$ and $M^{[2]}M$ have rank $80$.
Thus Lemmas~\ref{lem:hasse-line-certificate} and
\ref{lem:char2-defect-one-inertia} give transvection generation.
Apply Lemma~\ref{lem:primitive-prime-transvection-test}.

The complete matrices, factors, ranks and section coefficients
are in \path{outputs/product_degree3_certificates.json};
\path{code/verify_product_degree3.py} reconstructs them, independently
computes both $81\times81$ characteristic polynomials over the
integers before reduction modulo two, and repeats all the line
rank calculations.  Its exact output is
\path{outputs/verify_product_degree3.json}.
\end{proof}

\begin{lemma}[The required jet separation]
\label{lem:cubic-power-jets}
If $L$ is very ample on a smooth variety, $L^3$ separates second
jets at one point and first jets at two distinct points.
On a smooth projective surface its complete linear system has an irreducible
discriminant whose generic member has exactly one ordinary node;
the nonreduced locus has codimension at least two.
\end{lemma}

\begin{proof}
Ratios of sections of $L$ give local coordinates, so products of
at most two such ratios span the local ring modulo the cube of
the maximal ideal.  Multiplication by a section nonzero at the
point identifies this with second-jet generation by $L^3$.
For distinct points $P,Q$, take a section $m$ of $L$ vanishing
at $Q$ and nonzero at $P$.  Sections $m^2t$ with $t\in H^0(L)$
vanish to order at least two at $Q$ and give arbitrary first
jets at $P$.  Interchanging the points proves the assertion.

The incidence of a section with a specified singular point is
a projective bundle over the surface, since value and first
derivative impose three independent linear conditions.  It is
irreducible of dimension one less than the parameter space.
Second-jet separation makes an ordinary quadratic node a nonempty
open condition in this incidence, including in characteristic two.
Two specified distinct singular points impose six independent
conditions; varying the two points adds four dimensions.
Thus the locus with two distinct singular points has codimension
at least two in the parameter space.  The generic singular member
has exactly one ordinary node, and the discriminant is an
irreducible divisor.  A nonreduced divisor is singular along a
positive-dimensional locus, so belongs to the codimension-two
locus just described.
\end{proof}

\begin{theorem}[Transvection generation in every degree]
\label{thm:product-all-degree-transvections}
For the fixed $X=E\times E$ and $L$ of
Theorem~\ref{thm:product-char2-two-full-groups}, every geometric
monodromy group $G_{{\rm geom},d}$, $d\ge1$, is generated by
transvections.  Equivalently for the proof, every geometric
irreducible component of the full degree-$d$ Hasse divisor has
generic stable rank exactly $9d^2-2$ on the relative coherent
bundle of rank $9d^2-1$.
\end{theorem}

\begin{proof}
The divisor statement has been checked in degrees one, two and
three by the exact line certificates above.  Fix $d\ge4$, put
$e=d-3$, and work over the algebraic closure.
Choose a general smooth ordinary curve $D\in|L^e|$.
Such curves exist for every $e\ge1$ by
Theorem~\ref{thm:product-exact-fixed-covectors}.
Write $P_3=\mathbf P H^0(X,L^3)$ and let $C$ vary in this space.

We describe the bad loci for a general pencil in $P_3$.
The base Hasse divisor has generic stable rank loss one, and its
exceptional stable-rank locus has codimension at least two in
$P_3$.  By Lemma~\ref{lem:cubic-power-jets}, nonreduced curves and
curves worse than the generic one-node member form loci of
codimension at least two in the discriminant analysis relevant
here.  In particular every divisorial Hasse component has a
reduced generic member with at most one node.

For fixed smooth $D$, the incidence
\[
 \{(P,C): P\in D,\ s_C(P)=0,\ ds_C|_{T_PD}=0\}
\]
is a projective bundle over $D$ and is irreducible.  Its image
$\Delta_D\subset P_3$ is an irreducible divisor.  At its generic
point, $C$ is smooth and $C\cap D$ has one point of length two
and is otherwise reduced.  Indeed, contact of order at least
three imposes three conditions at one point of $D$, giving
codimension at least two after that point varies.  Two distinct
tangencies impose four independent conditions and the two points
vary in dimension two, again giving codimension at least two.
A singularity of $C$ at the tangency point adds a further linear
condition.  A singularity elsewhere, together with the tangency,
imposes five independent conditions at two points, which vary in
dimension three.  All these independence claims follow from
the two jet-separation assertions of
Lemma~\ref{lem:cubic-power-jets}.

For a general ordinary $D$, the generic $C$ on $\Delta_D$ is
ordinary.  Here is a dominance check to justify this assertion,
rather than a transversality assumption on Hasse strata.
Consider the joint incidence $(P,C,D)$ with both curves smooth
at $P$ and with equal tangent lines there.  Value and first-jet
generation for $L^3$ and $L^e$ identify it, locally over $P\in X$,
with the irreducible rank-one locus of two nonzero proportional
cotangent vectors, times affine kernel spaces.  Hence the joint
incidence is irreducible.  It dominates each of the two complete
linear systems: for a smooth curve in either system choose a
point on it and prescribe a proportional nonzero first jet for
the other section.  The inverse images of the two ordinary opens
are therefore nonempty dense opens of this same irreducible
incidence, so their intersection is dense.  Passing to a general
$D$ gives the asserted ordinary generic member of $\Delta_D$.
Consequently $\Delta_D$ is not a component of the base Hasse
divisor.  Their intersections and their exceptional loci have
codimension at least two in $P_3$.

The locus of $C$ containing $D$ is also of codimension at least
two.  For $e\le3$ it is the linear subspace
$s_D H^0(L^{3-e})$, whose codimension is respectively $45,72,80$;
for $e>3$ it is empty.  We may therefore choose a projective
pencil $\Lambda\subset P_3$ avoiding all the codimension-two
loci above.  Its point at infinity is chosen ordinary and
transverse to $D$.

For every member of this pencil, put $Y=C\cup D$.
There are no common components.  The exact sequence
\[
 0\longrightarrow\mathcal O_Y
 \longrightarrow\mathcal O_C\oplus\mathcal O_D
 \longrightarrow\mathcal O_{C\cap D}\longrightarrow0
\]
and exactness of stable parts give
\begin{equation}\label{eq:product-union-stable-rank}
 f(Y)=f(C)+f(D)+\#(C\cap D)(\bar k)-1,
\end{equation}
where $f$ denotes the stable dimension of $H^1(\mathcal O)$.
To see the finite-scheme term explicitly, Frobenius kills the
nilradical of $H^0(\mathcal O_{C\cap D})$ and is bijective on
its reduced geometric quotient.  Also $H^0(\mathcal O_C)=
H^0(\mathcal O_D)=\bar k$.  The coherent analogue uses the
scheme length instead of the number of points:
\[
 p_a(Y)=p_a(C)+p_a(D)+\operatorname{length}(C\cap D)-1.
\]
At a Hasse point of $C$, the intersection is reduced and $C$
loses exactly one stable dimension.  At a point of $\Delta_D$,
both components are ordinary and the number of intersection
points is one less than their scheme length.  These are distinct
sets of pencil parameters.  At every other parameter the union
is ordinary.  Thus every nonordinary union on the pencil has
stable rank loss exactly one.

Multiplication by the fixed defining section $s_D$ embeds this
pencil as a genuine projective line in $\mathbf P H^0(X,L^d)$.
Its point at infinity is ordinary, so every geometric component
of the full Hasse divisor meets the line at a finite parameter.
If such a component had generic stable rank loss at least two,
the closed stable-rank locus defined by the appropriate minors
of an iterated Frobenius matrix would contain the component and
all of those intersection points.  This contradicts
\eqref{eq:product-union-stable-rank}.  Hence every full-system
Hasse component has generic loss one.  The relative and absolute
losses agree, since $H^1(X,\mathcal O_X)$ has bijective Frobenius.
Lemma~\ref{lem:char2-defect-one-inertia} now proves transvection
generation in every degree.
\end{proof}

\begin{theorem}[All characteristic-two invariant covectors on products]
\label{thm:product-char2-all-parities-fixed-quotient}
Let $X=\prod_{i=1}^n E_i$ be a product of ordinary elliptic
curves in characteristic two, $n\ge2$, and
$L=\boxtimes_i\mathcal O_{E_i}(a_iO_i)$ with $a_i\ge3$.
For any $d\ge1$, let
$e=\#\{i: a_i d\text{ is even}\}$.
Then the entire geometric invariant covector space has dimension
\[
             \dim(V_d^\vee)^{G_{{\rm geom},d}}=2^e-1.
\]
After a finite extension splitting the ordinary Frobenius kernels
and normalizing the invariant top Cartier differential, the
arithmetic invariant space is the same space.
In particular, on the fixed $E\times E$ with $a_1=a_2=3$,
this dimension is zero for odd $d$ and three for even $d$.
\end{theorem}

\begin{proof}
Work over an algebraic closure and normalize the invariant top
form by Cartier.  The proof of
Theorem~\ref{thm:product-exact-fixed-covectors}, through
\eqref{eq:product-exact-Cartier-operator}, used no divisibility
assumption.  It identifies invariant covectors with exact lifts
$D\in\End(W)$, modulo $\mathbf F_2\operatorname{id}$, satisfying
\begin{equation}\label{eq:char2-projector-exact}
        \beta(s,D(s))=D(s)^{[2]}\quad\text{for every }s\in W,
\end{equation}
where $\beta$ is the symmetric relative Cartier pairing on
global sections and $\beta(s,s)=s^{[2]}$.

Every such $D$ is an idempotent.  An eigenvector with eigenvalue
$a$ in \eqref{eq:char2-projector-exact} gives $a=a^2$.
If a nontrivial Jordan block remained, subtract its eigenvalue
times the identity and choose $v,w\ne0$ with $D(v)=w$ and
$D(w)=0$.  Substituting $v+tw$ in the same equation makes the
coefficient of $t$ on the left equal to $\beta(w,w)=w^{[2]}$,
whereas the right side is independent of $t$, a contradiction.
Thus $W=A\oplus B$ is its $1$- and $0$-eigenspace decomposition.
Substituting $a+tb$ shows that
\begin{equation}\label{eq:char2-projector-orthogonality}
                    \beta(A,B)=0.
\end{equation}

Let $F:X\to X^{(2)}$ be relative Frobenius and
$\mathcal F=F_*L^d$.  Finite duality gives a perfect symmetric
pairing
\[
 \mathcal F\otimes\mathcal F\longrightarrow L^{d,(2)},
\]
obtained by multiplying sections and applying the normalized
Cartier trace.  Its pairing on global sections is $\beta$.
The bundle $\mathcal F$ is globally generated.  On a factor
$E_i$, evaluation on a Frobenius fibre is evaluation on a
length-two subscheme; $H^1(\mathcal O_{E_i}(a_idO_i)(-Z))=0$
because its degree is $a_i d-2>0$.  Tensoring these surjections
proves global generation on the product.

The projector $D$ preserves every evaluation kernel.
Indeed, write a global section whose value in a fibre is zero
as $a+b$ with $a\in A,b\in B$.  Its pairing with every section
is zero in that fibre.  By
\eqref{eq:char2-projector-orthogonality}, the value of $a$ pairs
to zero with $A$ and with $B$, hence with the whole fibre by
global generation.  Perfectness makes that value zero, and the
same holds for $b$.  The evaluation kernel is a vector bundle;
vanishing on all closed fibres therefore gives a descended
bundle endomorphism $\mathcal D\in\End(\mathcal F)$.
Its action on global sections is the original $D$.

We compute this endomorphism algebra directly.  For an ordinary
elliptic curve and a line bundle $N$ of odd degree $m$, the
rank-two bundle $F_*N$ has degree $m$ and is stable.
If a line subbundle $A\subset F_*N$ existed with degree at least
$m/2$, adjunction would give a nonzero map $F^*A\to N$ and
$2\deg A\le m$, impossible for the proposed integral degree.
Thus every line subbundle has degree strictly below the slope.
A noninvertible nonzero endomorphism would have a rank-one
kernel and image, both of degree below the slope (saturating the
image only raises its degree), contradicting additivity of
degree.  Over the algebraically closed field it follows that
\[
                     \End(F_*N)=\bar k.
\]

For even degree, $N=F^*N'$ and the ordinary splitting gives
$F_*N=N'\oplus(N'\otimes T)$, where $T$ is the nontrivial
two-torsion character line bundle.  The two summands are
nonisomorphic and have no homomorphisms between them.
This is the dimension-one case of the opened ordinary-Frobenius
splitting in \cite[Lemma~4.5, p.~11]{SannaiTanaka2016}.
Their cross Cartier pairing is zero, since the trace kills the
nontrivial character.

Relative Frobenius on the product is the product of the relative
Frobenii.  The bundle $\mathcal F$ is therefore a direct sum of
$2^e$ external products, one for each choice of a character on
the even-degree factors.  Each external product has endomorphism
algebra $\bar k$, by the degree-zero Kunneth formula.  Distinct
summands have no homomorphisms, because on one even-degree factor
the corresponding degree-zero nontrivial character line has no
sections.  Thus
\[
                      \End(\mathcal F)=\bar k^{\,2^e}.
\]
Every exact $D$ is consequently a sum of some of these $2^e$
projectors.  Conversely the summands are orthogonal for the
Cartier pairing; every such projector satisfies
\eqref{eq:char2-projector-exact}.  The exact-lift space therefore
has dimension $2^e$ over $\mathbf F_2$.
Quotienting by the identity gives the asserted $2^e-1$.
All splittings and the normalized top form descend after a
finite extension, which also proves the arithmetic statement.
\end{proof}

\begin{remark}[What the uniform results leave open]
For the fixed surface, we now know the exact generic rank, the
entire invariant quotient in every degree, and that geometric
monodromy is generated by transvections in every degree.
The complete images are identified above in degrees one, two
and three.  A transvection-generated group with the stated
invariant quotient can still have further invariant subspaces
or a proper translation subgroup.  The two uniform theorems do
not remove those possibilities for all $d$, and no all-degree
full-stabilizer assertion is made here.
\end{remark}
